Com o objetivo de expandir a presença humana além da Terra, outras missões estão em fase de pesquisa e desenvolvimento para possibilitar a colonização de outros planetas, como Marte.

\subsection{Mars Sample Return}

A missão \textit{Mars Sample Return} (MSR) \citep{nasa_mars_sample_return_2025} é uma iniciação conjunta da NASA e da Agência Espacial Europeia (ESA) que deseja coletar e trazer para a Terra amostras de rochas e solo marcianos para realizar análise de suas propriedades em laboratórios terrestres. Com o intuito de preparar o caminho para futuros exploradores humanos em Marte e de aprofundar a compreensão sobre Marte e o sistema solar, essa campanha com diversas missões derivadas dela, é considerada uma das mais desafiadoras já planejadas.

O conceito da missão envolve várias etapas e componentes, de início, o rover \textit{Perseverance}, que já está em Marte, coleta e armazena amostras de rochas e solo e os armazenam em tubos selados. Depois disso, uma futura missão enviará o \textit{Sample Retrieval Lander}, que pousará próximo ao \textit{Perseverance}, na cratera Jezero, para receber as amostras coletadas. Este módulo de aterrissagem contará com um pequeno foguete, o \textit{Mars Ascent Vehicle} \citep{nasa_mars_ascent_vehicle_2025}, no qual as amostras serão carregadas para serem lançadas na órbita marciana. Uma vez em órbita, os \textit{Mars Ascent Vehicle} (e as amostras) serão capturadas pelo \textit{Earth Return Orbiter} \citep{nasa_earth_return_orbiter_2025}, uma espaçonave projetada para trazê-las de volta à Terra para análises aprofundadas.

A coleta e o retorno dessas amostras marcianas são etapas significativas para responder a questões fundamentais sobre a possibilidade de vida em Marte, a história geológica do planeta e sua evolução climática. As amostras retornadas permitirão realizar estudos para responder a dúvidas e também, possibilitando descobertas que não seriam viáveis apenas com os dados coletados por sondas e \textit{rovers} em Marte.

Além disso, existe a exploração de asteroides, onde sistemas de navegação autônomos realizam aproximações com corpos irregulares e de baixa gravidade.

\subsection{Navegação autônoma na proximidade de asteroides}

A navegação autônoma de espaçonaves na proximidade de asteroides próximos à Terra (em inglês, \textit{Near Earth Asteroid} ou NEOs) possui complexidades adicionais envolvidas na estimativa precisa de suas trajetórias e na necessidade de operações de alta precisão. Para contornar essas dificuldades extras, se faz necessário a utilização de múltiplos sensores embarcados, como câmeras, sensores de atitude e LiDAR, e tem sido explorado para melhorar a navegação desses veículos espaciais e permitir a operação coordenada de formações de espaçonaves. Conforme explica \citep{vetrisano2015autonomous}, a fusão de dados de medições interespaciais com as estimativas calculadas de posição relativas ao asteroide pode melhorar a robustez do sistema imprecisões nos dados e até falhas.

O uso de medições relativas entre as espaçonaves, durante um voo de formação, pode contribuir para a mitigação de erros associados a modelos de perturbações gravitacionais incertos e variações no formato e rotação dos asteroides. As estratégias que combinam medições de linha de visada com sensores baseados no efeito Doppler do Sol, conforme explica \citep{yim2000autonomous}, podem ser utilizadas para determinar a órbita do asteroide com maior precisão.

O voo em formações de espaçonaves que orbitam um asteroide requer controle autônomo devido à natureza não linear, e imprevisível, da dinâmica orbital dos asteroides, além disso, a radiação solar pode influenciar negativamente na estabilidade da órbita, e irá exigir técnicas mais avançadas de controle para manter a formação desejada. Métodos de controle baseados em Lyapunov foram propostos para garantir que as espaçonaves permaneçam próximas a uma trajetória de referência \citep{vetrisano2015autonomous}.


